\documentclass[10pt,aspectratio=169]{beamer}
\usetheme{Madrid}
\usecolortheme{default}

\usepackage[utf8]{inputenc}
\usepackage[brazil]{babel}
\usepackage{listings}
\usepackage{xcolor}
\usepackage{float}

% =========================================================
% PALETA DE CORES - EQUIPE IMPERADOR
% =========================================================
\definecolor{imperadorblue}{RGB}{8, 24, 50}      % Azul escuro (Fundo e painéis)
\definecolor{imperadoryellow}{RGB}{244, 195, 0}  % Amarelo vibrante (Destaques e títulos)
\definecolor{codebg}{RGB}{240, 242, 245}         % Fundo dos blocos de código

% Configuração das cores do Beamer
\setbeamercolor{palette primary}{bg=imperadorblue,fg=white}
\setbeamercolor{palette secondary}{bg=imperadoryellow,fg=imperadorblue}
\setbeamercolor{palette tertiary}{bg=imperadorblue,fg=white}
\setbeamercolor{titlelike}{parent=palette primary,fg=imperadoryellow}
\setbeamercolor{frametitle}{bg=imperadorblue,fg=imperadoryellow}
\setbeamercolor{structure}{fg=imperadorblue}
\setbeamercolor{item}{fg=imperadoryellow}
\setbeamercolor{block title}{bg=imperadoryellow,fg=imperadorblue}
\setbeamercolor{block body}{bg=imperadorblue!5,fg=black}

% Configuração para estilização dos blocos de código (Terminal)
\lstset{
    backgroundcolor=\color{codebg},
    basicstyle=\ttfamily\small,
    breaklines=true,
    keywordstyle=\color{imperadorblue}\bfseries,
    stringstyle=\color{imperadorblue!80},
    frame=single,
    rulecolor=\color{imperadorblue!30},
    frameround=tttt,
    showstringspaces=false
}

% =========================================================
% METADADOS DA APRESENTAÇÃO
% =========================================================
\title[Curso de Git - Imperador]{Curso prático de Git}
\subtitle{Do Básico ao Avançado}
\author{Lambari}
\institute{Equipe Imperador - UTFPR Curitiba}
\date{Julho 2026}

\begin{document}

% Slide de Título
\begin{frame}
    \titlepage
\end{frame}

% Slide de Agenda
\begin{frame}{Tópicos}
    \tableofcontents
\end{frame}

% =========================================================
% MÓDULO 1
% =========================================================
\section{Módulo 1: Básico (Github, repositório, etc.)}

\begin{frame}{Módulo 1: Básico (Github, repositório, etc.)}
    \begin{block}{Objetivo}
        Compreender a filosofia do Git, configurar o ambiente local e dominar o fluxo de trabalho essencial para arquivos individuais.
    \end{block}
    \vspace{0.5cm}
    \textbf{Tópicos abordados:}
    \begin{itemize}
        \item O que é o Git vs. GitHub
        \item Configuração inicial do ambiente
        \item Ciclo de vida de um arquivo
        \item Comandos fundamentais locais e remotos
    \end{itemize}
\end{frame}

\begin{frame}[fragile]{Git vs. GitHub}
    \begin{columns}
        \begin{column}{0.5\textwidth}
            \textbf{Git}
            \begin{itemize}
                \item Sistema de controle de versão \textbf{local} e distribuído.
                \item Funciona inteiramente no terminal do computador.
                \item Tira "fotos" (commits) do histórico do código.
            \end{itemize}
        \end{column}
        \begin{column}{0.5\textwidth}
            \textbf{GitHub / GitLab}
            \begin{itemize}
                \item Plataforma de hospedagem \textbf{na nuvem}.
                \item Interface gráfica para gerenciar repositórios.
                \item Adiciona ferramentas de colaboração (Pull Requests).
            \end{itemize}
        \end{column}
    \end{columns}
\end{frame}

\begin{frame}[fragile]{Configuração Inicial do Ambiente}
    Antes do primeiro commit, o Git precisa saber quem é você para registrar a autoria corretamente.
    \vspace{0.3cm}
    \begin{lstlisting}[language=bash]
# Definir seu nome globalmente
$ git config --global user.name "Seu Nome Completo"

# Definir seu e-mail de trabalho
$ git config --global user.email "seu_email@imperador.com"

# Verificar configuracoes ativas
$ git config --list
    \end{lstlisting}
\end{frame}

\begin{frame}{O Ciclo de Vida dos Arquivos}
    No Git, os arquivos passam por estados específicos antes de serem consolidados:
    \vspace{0.3cm}
    \begin{description}
        \item[\textcolor{imperadorblue}{\textbf{Untracked}}] Arquivo novo no diretório, o Git ainda não o monitora.
        \item[\textcolor{imperadorblue}{\textbf{Modified}}] Arquivo monitorado que sofreu alterações locais.
        \item[\textcolor{imperadorblue}{\textbf{Staged}}] Arquivo selecionado e pronto para fazer parte do próximo commit.
        \item[\textcolor{imperadorblue}{\textbf{Committed}}] Alterações salvas com segurança no banco de dados local.
    \end{description}
\end{frame}

\begin{frame}[fragile]{Iniciando um Repositório: git init}
    O comando \texttt{git init} é utilizado para criar um novo repositório Git vazio ou reinicializar um já existente.
    \vspace{0.3cm}
    
    \textbf{Atenção:} Na prática do dia a dia, é mais comum criarmos o repositório diretamente em uma plataforma como o \textbf{GitHub} e depois clonarmos para a nossa máquina, tornando o \texttt{git init} um comando raramente utilizado manualmente pela nossa equipe.
    
    \vspace{0.3cm}
    \begin{lstlisting}[language=bash]
# Inicializar um repositorio vazio na pasta atual
$ git init
# Clonar um repositorio da nuvem na past atual
$ git clone <URL>
    \end{lstlisting}
\end{frame}

\begin{frame}[fragile]{Verificando o Estado: git status}
    O comando \texttt{git status} exibe o estado do diretório de trabalho e da área de preparação (staging area).
    \vspace{0.3cm}
    
    Ele permite ver quais arquivos foram modificados, quais estão prontos para o próximo commit e quais ainda não estão sendo rastreados pelo Git. É o seu "radar" de desenvolvimento.
    \vspace{0.3cm}
    \begin{lstlisting}[language=bash]
# Verificar o estado atual dos arquivos e modificacoes
$ git status
    \end{lstlisting}
    \begin{figure}[H]
        \centering
        \includegraphics[width=0.7\linewidth]{GitStatus.png}
        \label{fig:enter-label}
    \end{figure}
\end{frame}

\begin{frame}[fragile]{Preparando Arquivos: git add}
    O comando \texttt{git add} adiciona uma alteração do seu diretório de trabalho à área de preparação (staging area).
    \vspace{0.3cm}
    
    Isso avisa ao Git que você quer incluir as atualizações de um arquivo específico no próximo commit.
    \vspace{0.3cm}
    \begin{lstlisting}[language=bash]
# Adicionar o script modificado para a area de preparacao
$ git add sensor_velocidade.cpp
# Adiciona todas as modificações para a area de preparacao
$ git add .
    \end{lstlisting}
\end{frame}

\begin{frame}[fragile]{Salvando as Alterações: git commit}
    O comando \texttt{git commit} captura um "instantâneo" das alterações que estão atualmente na área de preparação (staging area).
    \vspace{0.3cm}
    
    Todo commit deve ser acompanhado de uma mensagem clara e descritiva para explicar o motivo da alteração no histórico do projeto.
    \vspace{0.3cm}
    \begin{lstlisting}[language=bash]
# Salvar permanentemente o estado atual com uma breve ferramentas
$ git commit -m "feat(sensores): adiciona leitura de pulsos do sensor Hall"
# Salvar permanentemente o estado atual com uma mensagem mais detalhada em um editor de texto
$ git commit 
    \end{lstlisting}
\end{frame}

\begin{frame}[fragile]{Trabalhando com o Repositório Remoto}
    Sincronização com o servidor em nuvem da equipe:
    \vspace{0.2cm}
    \begin{lstlisting}[language=bash]
# Clonar um repositorio existente para sua maquina local
$ git clone <URL>

# Enviar seus commits locais para a nuvem
$ git push origin main 

# Buscar as novidades da nuvem e atualizar seu codigo local
$ git pull origin main
    \end{lstlisting}
\end{frame}

% =========================================================
% MÓDULO 2
% =========================================================
\section{Módulo 2: Intermediário (Branches, Merges e .gitignore)}

\begin{frame}{Módulo 2: Intermediário (Branches, Merges e .gitignore)}
    \begin{block}{Objetivo}
        Aprender a trabalhar em paralelo usando ramificações, integrar desenvolvimentos diferentes e gerenciar arquivos descartáveis.
    \end{block}
    \vspace{0.5cm}
    \textbf{Tópicos abordados:}
    \begin{itemize}
        \item Conceito e comandos de Branches
        \item Integração de código com Merge
        \item Resolução controlada de conflitos
        \item Configuração estratégica do \texttt{.gitignore}
    \end{itemize}
\end{frame}

\begin{frame}{O Propósito das Branches}
    Uma \textbf{branch} (ramificação) é uma linha de desenvolvimento independente. Funciona como um "universo paralelo" onde você pode alterar arquivos livremente sem afetar o código oficial (a branch \texttt{main}).
    \vspace{0.3cm}
    
    \textbf{Por que usar?}
    \begin{itemize}
        \item \textbf{Isolamento:} Evita que seu código em desenvolvimento quebre a versão principal.
        \item \textbf{Trabalho Paralelo:} Vários membros da equipe podem programar simultaneamente.
        \item \textbf{Organização:} Permite revisar a funcionalidade antes de integrá-la.
    \end{itemize}
    \vspace{0.2cm}
    
    \begin{figure}[H]
        \centering
        \includegraphics[width=0.6\linewidth]{GitBranches.png}
    \end{figure}
\end{frame}

\begin{frame}[fragile]{Ramificações (Branches)}
    Uma branch isola o seu desenvolvimento. 
    
    \textbf{Regra na equipe:} Nunca faça commits diretamente na branch \texttt{main}.
    \vspace{0.2cm}
    \begin{lstlisting}[language=bash]
# Lista as branches que voce ja utilizou nesse repositorio
$ git branch 

# Deletar uma branch do seu computador
$ git branch -D <nome da branch>

# Alterar para uma branch
$ git switch feature/filtros-sinais

# Comando alternativo (mais antigo) para alterar de branch
$ git checkout feature/filtros-sinais
    \end{lstlisting}
\end{frame}

\begin{frame}[fragile]{Fusão de Código: Unindo Forças (Merge)}
    Como integrar o trabalho finalizado de volta à linha principal do projeto:
    \vspace{0.2cm}
    \begin{lstlisting}[language=bash]
# 1. Voltar para a branch principal
$ git switch main

# 2. Atualizar a main local com o servidor
$ git pull origin main

# 3. Fundir branches (feito no PR do GitHub)
$ git merge feature/filtros-sinais

# 4. Remover a branch local que ja foi integrada
$ git branch -d feature/filtros-sinais
    \end{lstlisting}
\end{frame}

\begin{frame}[fragile]{Lidando com Conflitos de Frente}
    Conflitos ocorrem quando duas pessoas alteram a mesma linha do mesmo arquivo. O Git pausa o processo e exibe os marcadores:
    \vspace{0.2cm}
    \begin{lstlisting}[language=python]
<<<<<<< HEAD
limite_rpm = 3800  # Codigo atual na Main
=======
limite_rpm = 4000  # Alteracao feita na sua Branch
>>>>>>> feature/ajuste-motor
    \end{lstlisting}
    \vspace{0.1cm}
    \textbf{Resolução:} Abra o arquivo, analise com a equipe qual lógica deve prevalecer, remova os marcadores (\texttt{<<<}, \texttt{===}, \texttt{>>>}), salve, execute \texttt{git add} e finalize com \texttt{git commit}.
\end{frame}

\begin{frame}[fragile]{O Arquivo .gitignore}
    Evite que arquivos gerados por compilação, logs temporários ou pastas de IDEs sujem o histórico do repositório.
    \vspace{0.2cm}
    Basta criar um arquivo chamado \texttt{.gitignore} na raiz com as extensões ou pastas mapeadas:
    \vspace{0.2cm}
    \begin{lstlisting}[language=bash]
# Ignorar binarios e objetos de compilacao
*.o
*.bin
*.elf
build/
.vscode/
/logs/*.csv
    \end{lstlisting}
\end{frame}

% =========================================================
% MÓDULO 3
% =========================================================
\section{Módulo 3: Avançado (Revisão e Git Rebase)}

\begin{frame}{Módulo 3: Avançado (Revisão e Git Rebase)}
    \begin{block}{Objetivo}
        Dominar ferramentas avançadas para manipulação fina do histórico, limpeza de árvores de commits e recuperação de dados perdidos.
    \end{block}
    \vspace{0.5cm}
    \textbf{Tópicos abordados:}
    \begin{itemize}
        \item Armazenamento temporário com \texttt{git stash}
        \item Organização linear com \texttt{git rebase}
        \item Importação cirúrgica com \texttt{git cherry-pick}
        \item Recuperação extrema com \texttt{git reflog}
    \end{itemize}
\end{frame}

\begin{frame}[fragile]{Guardando Trabalho Incompleto: git stash}
    Se você estiver trabalhando em um código complexo e surgir um bug urgente na pista que exige troca imediata de branch:
    \vspace{0.2cm}
    \begin{lstlisting}[language=bash]
# Salva o progresso inacabado em um compartimento isolado e limpa o terminal
$ git stash save -m "Progressodo desenvolvimento da telemetria"

# Agora voce esta livre para alternar para a main e corrigir o problema
$ git switch main
# ... correcao do bug ...

# De volta à sua branch de trabalho, resgate as modificacoes salvas
$ git switch feature/telemetria
$ git stash pop
    \end{lstlisting}
\end{frame}

\begin{frame}[fragile]{Histórico Linear: Rebase vs. Merge}
    Enquanto o \texttt{merge} cria um commit de junção explícito, o \texttt{rebase} reposiciona os commits da sua branch diretamente no topo da branch de destino, gerando um histórico limpo e sem ramificações visuais complexas.
    \vspace{0.2cm}
    \begin{lstlisting}[language=bash]
# Atualiza sua branch aplicando suas modificacoes sobre a main atualizada
$ git switch feature/minha-feature
$ git rebase main
    \end{lstlisting}
    \vspace{0.1cm}
    \textbf{Aviso importante:} Nunca utilize rebase em branches públicas/compartilhadas com outros membros da equipe para evitar dessincronização de históricos.
\end{frame}

\begin{frame}[fragile]{A Caixa-Preta do Git: git reflog}
    Se uma branch foi deletada acidentalmente ou um comando de reset removeu arquivos válidos, o \texttt{reflog} salva o dia registrando cada alteração local do ponteiro HEAD.
    \vspace{0.2cm}
    \begin{lstlisting}[language=bash]
# Exibir o historico completo de operacoes locais
$ git reflog

# Saida esperada no terminal:
# e45a2d1 HEAD@{0}: reset: moving to HEAD~1
# a78b112 HEAD@{1}: commit: Adiciona calculo de G-Force

# Retornar o estado do projeto exatamente para o instante anterior ao erro:
$ git reset --hard a78b112
    \end{lstlisting}
\end{frame}

% =========================================================
% MÓDULO 4
% =========================================================
\section{Módulo 4: O Processo da Equipe (Implementação)}

\begin{frame}{Módulo 4: O Processo da Equipe (Implementação)}
    \begin{block}{Objetivo}
        Estabelecer a metodologia de trabalho da equipe. Padronizar a integração de código para garantir que o repositório sempre esteja funcional e validado.
    \end{block}
    \vspace{0.5cm}
    \textbf{Tópicos abordados:}
    \begin{itemize}
        \item O modelo GitHub Flow adaptado ao ambiente competitivo
        \item Pull Requests (PRs) e Code Review
        \item Padronização com Commits Semânticos
    \end{itemize}
\end{frame}

\begin{frame}{Fluxo Operacional: GitHub Flow}
    Para a dinâmica veloz, adota-se o formato simplificado:
    \vspace{0.3cm}
    \begin{enumerate}
        \item A branch \texttt{main} reflete o software funcional e testado no protótipo.
        \item Desenvolvimentos ocorrem em branches nomeadas conforme o escopo (ex: \texttt{feature/sensor-freio}).
        \item Commits são enviados para o servidor remoto.
        \item A integração final exige a abertura de um \textbf{Pull Request} (PR).
    \end{enumerate}
\end{frame}

\begin{frame}{Padronização de Branches (Parte 1)}
    Para manter a organização do repositório, adotamos uma convenção para as branches principais:
    \vspace{0.3cm}
    \begin{itemize}
        \item \textbf{\texttt{main}}: Contém o código testado e 100\% funcional, pronto para ir para o protótipo.
        \item \textbf{\texttt{main/YYYY-MM-DD}}: Branches de "snapshot". Salvam o estado exato do código em dias críticos (ex: competições), como \texttt{main/2026-07-05}.
    \end{itemize}
\end{frame}

\begin{frame}{Padronização de Branches (Parte 2)}
    As branches temporárias (de trabalho diário) seguem prefixos específicos que indicam o seu propósito:
    \vspace{0.3cm}
    \begin{itemize}
        \item \textbf{\texttt{feature/}}: Utilizadas para desenvolvimentos de funções (ex: \texttt{feature/leitura-gps}).
        \item \textbf{\texttt{bug/}}: Utilizadas exclusivamente para correção de bugs no código (ex: \texttt{bug/falha-telemetria}).
        \item \textbf{\texttt{test/}}: Branches para desenvolver códigos voltados para testes (ex: \texttt{test/validacao-inversor}).
        \item \textbf{\texttt{doc/}}: Utilizadas exclusivamente para atualização de documentação em uma branch.
    \end{itemize}
\end{frame}

\begin{frame}{Segurança do Código: Pull Requests e Code Review}
    \begin{itemize}
        \item \textbf{Branches Protegidas:} O envio direto de commits para a branch \texttt{main} deve ser bloqueado nas configurações do servidor remoto.
        \item \textbf{Revisão Técnica obrigatória:} Todo Pull Request precisa ser revisado e validado por pelo menos um outro integrante da área de eletrônica antes de ser integrado.
        \item \textbf{Impacto Direto:}
        \begin{itemize}
            \item Previne que falhas de sintaxe cheguem ao carro antes dos testes de campo.
            \item Garante que mais de um membro da equipe domine o funcionamento daquela seção do software.
        \end{itemize}
    \end{itemize}
\end{frame}

\begin{frame}[fragile]{Mensagens de Commit Padronizadas}
    Evite mensagens genéricas como "ajustes" ou "agora vai". Adote a convenção de Commits Semânticos:
    \vspace{0.2cm}
    \begin{lstlisting}[language=bash]
# Estrutura padrao: tipo(subsistema): descricao em letras minusculas
$ git commit -m "feat(telemetria): adiciona transmissao via modulo LoRa"
$ git commit -m "fix(dashboard): corrige escala do grafico de RPM"
$ git commit -m "docs(readme): atualiza esquema de pinagem do microcontrolador"
$ git commit -m "test(pwt): otimizacao cvt"
    \end{lstlisting}
    \vspace{0.1cm}
    \textbf{Tipos recomendados:}
    \begin{itemize}
        \item \textbf{feat}: Implementação de uma nova funcionalidade ou bloco de código.
        \item \textbf{fix}: Correção de uma falha ou bug operacional.
        \item \textbf{docs}: Alterações estritas na documentação técnica ou comentários.
        \item \textbf{test}: Código de um teste específico para um ou mais subsistemas.
    \end{itemize}
\end{frame}

\begin{frame}[fragile]{Passo a Passo do Desenvolvimento}
    Fluxo de trabalho ideal do início ao fim de uma nova tarefa:
    \vspace{0.2cm}
    \begin{enumerate}
        \item \textbf{Criar a branch no GitHub} (ex: \texttt{feature/leitura-gps})
        \item \textbf{Atualizar repositório local:} \texttt{git pull}
        \item \textbf{Alternar para a branch:} \texttt{git switch feature/leitura-gps}
        \item \textbf{Desenvolver:} Escrever o código, testar, etc.
        \item \textbf{Atualizar branch (se a main mudou):} \texttt{git rebase main} (resolver eventuais conflitos)
        \item \textbf{Preparar alterações:} \texttt{git add .} (ou arquivo específico)
        \item \textbf{Salvar:} \texttt{git commit -m "feat(telemetria): rotina do GPS"}
        \item \textbf{Enviar:} \texttt{git push origin feature/leitura-gps}
        \item \textbf{Abrir o PR} (Pull Request) no GitHub.
        \item \textbf{Realizar o Merge} (após aprovação no Code Review).
    \end{enumerate}
\end{frame}

\end{document}